\section{Diagramas de Classe}

O sistema segue uma arquitetura de módulos, cada módulo possui seu próprio domínio, casos de uso e camada de infraestrutura. Cada diagrama exibe as três camadas da arquitetura do módulo: \textbf{Domain} (entidades, objetos de valor e enumerações), \textbf{Application} (casos de uso e portas) e \textbf{Infrastructure} (controladores, repositórios e adaptadores).

\subsection{Usuário}

Na Figura~\ref{fig:classe-user} são apresentadas as classes do módulo de gerenciamento de usuários.

\begin{figure}[H]
    \caption{Diagrama de classe do usuário}
    \label{fig:classe-user}
    \centering
    \includegraphics[width=0.95\textwidth,height=0.75\textheight,keepaspectratio]{../../../out/src/plantuml/classe/simples/user_simples.png}

    {\footnotesize \textbf{Fonte:} O autor.}
\end{figure}

\subsection{Sala}

Na Figura~\ref{fig:classe-building}, Figura~\ref{fig:classe-room-type}, Figura~\ref{fig:classe-resource} e Figura~\ref{fig:classe-room} são apresentadas respectivamente as classes de \texttt{Building}, \texttt{RoomType}, \texttt{Resource} e \texttt{Room}, que compõem o módulo de gerenciamento de salas.

\begin{figure}[H]
    \caption{Diagrama de classe do prédio}
    \label{fig:classe-building}
    \centering
    \includegraphics[width=0.95\textwidth,height=0.75\textheight,keepaspectratio]{../../../out/src/plantuml/classe/simples/room/room_building.png}

    {\footnotesize \textbf{Fonte:} O autor.}
\end{figure}

\begin{figure}[H]
    \caption{Diagrama de classe do tipo de sala}
    \label{fig:classe-room-type}
    \centering
    \includegraphics[width=0.95\textwidth,height=0.75\textheight,keepaspectratio]{../../../out/src/plantuml/classe/simples/room/room_room_type.png}

    {\footnotesize \textbf{Fonte:} O autor.}
\end{figure}

\begin{figure}[H]
    \caption{Diagrama de classe do recurso}
    \label{fig:classe-resource}
    \centering
    \includegraphics[width=0.95\textwidth,height=0.75\textheight,keepaspectratio]{../../../out/src/plantuml/classe/simples/room/room_resource.png}

    {\footnotesize \textbf{Fonte:} O autor.}
\end{figure}

\begin{figure}[H]
    \caption{Diagrama de classe da sala}
    \label{fig:classe-room}
    \centering
    \includegraphics[width=0.95\textwidth,height=0.75\textheight,keepaspectratio]{../../../out/src/plantuml/classe/simples/room/room_room.png}

    {\footnotesize \textbf{Fonte:} O autor.}
\end{figure}

\subsection{Evento}

Na Figura~\ref{fig:classe-event} e Figura~\ref{fig:classe-period} são apresentadas respectivamente as classes do evento com sua solicitação de alteração e as classes do período, ambas do módulo de gerenciamento de eventos.

\begin{figure}[H]
    \caption{Diagrama de classe do evento e da solicitação}
    \label{fig:classe-event}
    \centering
    \includegraphics[width=0.95\textwidth,height=0.75\textheight,keepaspectratio]{../../../out/src/plantuml/classe/simples/event/event_and_request.png}

    {\footnotesize \textbf{Fonte:} O autor.}
\end{figure}

\begin{figure}[H]
    \caption{Diagrama de classe do período}
    \label{fig:classe-period}
    \centering
    \includegraphics[width=0.95\textwidth,height=0.75\textheight,keepaspectratio]{../../../out/src/plantuml/classe/simples/event/event_period.png}

    {\footnotesize \textbf{Fonte:} O autor.}
\end{figure}

\subsection{Relatório}

Na Figura~\ref{fig:classe-report} são apresentadas as classes do módulo de relatórios.

\begin{figure}[H]
    \caption{Diagrama de classe do relatório}
    \label{fig:classe-report}
    \centering
    \includegraphics[width=0.95\textwidth,height=0.75\textheight,keepaspectratio]{../../../out/src/plantuml/classe/simples/report_simples.png}

    {\footnotesize \textbf{Fonte:} O autor.}
\end{figure}
